\documentclass[10pt,a4paper]{amsart}
\usepackage[margin=19mm]{geometry}
\usepackage{amsmath,amssymb,mathtools}
\usepackage[scr=rsfs,cal=euler]{mathalfa}
\usepackage{xfrac}
\usepackage[unicode,hidelinks,bookmarks=false]{hyperref}
\newcommand{\ie}{i.e.\ }
\newcommand{\cf}{cf.\ }
\newcommand{\resp}{respectively\ }
\newcommand{\Subsection}{\subsection}
\providecommand{\nobreakdash}{\nobreak\hbox{-}}
\setlength{\emergencystretch}{2em}
\multlinegap0pt
\usepackage{amssymb}
\usepackage[shortlabels]{enumitem}
\usepackage{tcolorbox}
\newtheorem{lemma}{Lemma}
\def\slashiii#1{\setbox0=\hbox{$#1$}#1\hskip-\wd0\hbox to\wd0{\hss\sl/\/\hss}}
\title{On the super-Liouville equations on the sphere}
\author{Mingyang Han, Chunqin Zhou}
\date{Source: arXiv:2509.16712v5 (CC BY 4.0)}
\begin{document}
\maketitle

\noindent\emph{Source and licence.} On the super-Liouville equations on the sphere. Version arXiv:2509.16712v5 (CC BY 4.0), distributed by arXiv under the Creative Commons Attribution 4.0 International licence, http://creativecommons.org/licenses/by/4.0/, as recorded in the arXiv licence metadata for this exact version and shown on the abstract page https://arxiv.org/abs/2509.16712v5. Full licence notice: \texttt{licenses/arxiv-2509.16712-CC-BY-4.0.txt}.\par\smallskip

\noindent\textit{Editorial scope.} Counted unit: the article's lemma numerator (source0.tex lines 2098--2123). The article's complete original proof of it is delivered with it (source0.tex lines 2125--2436, one \texttt{proof} environment of 312 lines). No mathematics is written, repaired or re-derived: the statement and the proof are the article's own lines, in the article's own notation, and no retained line is substituted or re-flowed.  No further source-local context is needed: every cross-reference of every retained line resolves inside the two delivered spans.  Nothing was omitted from any retained span, so the package carries no omission record.  No retained line contains a citation, so the wrapper carries no bibliography and no bibliography entry.  No retained line contains a figure environment, an image-inclusion command or drawing code, so no image is delivered and the package declares no asset dependency.  Editorial formatting only, in addition: an AMS \texttt{amsart} preamble that restates the article's own declarations and macros used by the retained lines, namely the environments \texttt{lemma} on separate counters, together with the standard packages those lines need.  The retained environments print in this document as Lemma 1 in order of appearance, whereas the article prints the same items with the numbering of its own full document.  The complete native source of the article as distributed in the arXiv e-print for this exact version is bundled unchanged as \texttt{sources/185397/source0.tex}.
\begin{lemma}
	Let \(p\in\mathbb S^n\), and choose geodesic normal coordinates centered at
	\(p\). After fixing a smooth local orthonormal spin frame, we identify the
	spinor bundle over the coordinate ball with the trivial bundle
	\(B_\rho(0)\times\mathbb C^N\). Thus a compactly supported spinor in the
	coordinate ball may be viewed as a \(\mathbb C^N\)-valued function. For \(\Phi\in C_c^\infty(B_R(0),\mathbb C^N)\), define
	\[
	\Phi_\varepsilon(\exp_p(\varepsilon y))
	=
	\varepsilon^{-\frac{n-1}{2}}\Phi(y),
	\]
	and set \(\Phi_\varepsilon=0\) outside \(B_{\varepsilon R}(p)\). Then
	\begin{equation}\label{numerator}
	\lim_{\varepsilon\to0}
	\int_{\mathbb S^n}
	\langle |\slashiii{D}|\Phi_\varepsilon,\Phi_\varepsilon\rangle\,dv_{g_0}
	=
	\int_{\mathbb R^n}
	\langle (-\Delta_{\mathbb R^n})^{1/2}\Phi,\Phi\rangle_{\mathbb C^N}\,dy,
	\end{equation}
	where \((-\Delta_{\mathbb R^n})^{1/2}\) denotes the Fourier multiplier
	\[
	\widehat{(-\Delta_{\mathbb R^n})^{1/2}\Phi}(\xi)	=|\xi|\widehat{\Phi}(\xi),
	\]
	acting componentwise on \(\mathbb C^N\)-valued functions.
\end{lemma}

\begin{proof}
	Fix $\chi \in C_c^\infty(B_\rho(p))$ with $\chi \equiv 1$ on
	$B_{\rho/2}(p)$ and $0 \le \chi \le 1$. For all $\varepsilon > 0$ small
	enough that $\varepsilon R < \rho/2$, the support of $\Phi_\varepsilon$ is
	contained in $B_{\rho/2}(p)$, so $\Phi_\varepsilon = \chi\Phi_\varepsilon$.
	Using this and the self-adjointness of $\chi$ as a multiplication operator
	($\langle u,\chi v\rangle = \langle \chi u, v\rangle$), we compute
	\[
	\begin{aligned}
		\langle |\slashiii{D}|\Phi_\varepsilon,\Phi_\varepsilon\rangle_{L^2(\mathbb{S}^n)}
		&= \langle |\slashiii{D}|(\chi\Phi_\varepsilon),\chi\Phi_\varepsilon
		\rangle_{L^2(\mathbb{S}^n)}\\
		&= \langle \chi\,|\slashiii{D}|\,\chi\,\Phi_\varepsilon,\Phi_\varepsilon
		\rangle_{L^2(\mathbb{S}^n)}.
	\end{aligned}
	\]
	Set $T := \chi\,|\slashiii{D}|\,\chi$. In the chosen coordinate chart and spin
	trivialization, $T$ may be represented as a properly supported classical
	pseudodifferential operator of order one on $\mathbb{R}^n$, after an
	arbitrary smooth extension of its local symbol outside the chart. This
	extension does not affect the quadratic form evaluated on $\Phi_\varepsilon$
	for sufficiently small $\varepsilon$, since $\operatorname{supp}\Phi_\varepsilon$
	lies strictly inside the chart. The Schwartz kernel of $T$ equals
	$K_T(x,z) = \chi(x)\,K_{|\slashiii{D}|}(x,z)\,\chi(z)$, which is compactly supported in
	$(x,z)$ but retains the standard conormal singularity along the diagonal
	$\{x = z\}$. The Kohn--Nirenberg symbol of \(T\) admits a classical asymptotic
	expansion. Fix an integer \(N_0\ge n+2\). We write it in the finite form
	\[
	a(x,\xi)
	=
	a_1(x,\xi)+a_0(x,\xi)+a_{-1}(x,\xi)+\cdots+a_{-N_0}(x,\xi)
	+r_{-N_0-1}(x,\xi),
	\]
	where \(a_j\in S^j\) and \(r_{-N_0-1}\in S^{-N_0-1}\). The principal symbol is
	\[
	a_1(x,\xi)=|\xi|_{g(x)}\,\mathrm{Id}_{\mathbb C^N},
	\qquad
	|\xi|_{g(x)}=\sqrt{g^{ij}(x)\xi_i\xi_j}.
	\]
	As usual, this identity for the principal symbol is understood on the
	high-frequency region \(|\xi|\ge c>0\); the low-frequency part is treated
	separately below by the fixed cutoff \(\varphi\). The lower-order terms \(a_j\), \(0\ge j\ge -N_0\), encode curvature,
	spin-connection, coordinate, and quantization contributions. The remainder
	\(r_{-N_0-1}\) is of sufficiently negative order.
	
	Representing $T$ by its Kohn--Nirenberg symbol up to a smoothing operator
	$R_\infty \in \Psi^{-\infty}$:
	\begin{equation}\label{eq:osc}
		\langle T\Phi_\varepsilon,\Phi_\varepsilon\rangle_{L^2}
		= (2\pi)^{-n}
		\iiint
		e^{i(x-z)\cdot\xi}\,
		\bigl\langle a(x,\xi)\Phi_\varepsilon(z),\Phi_\varepsilon(x)\bigr\rangle
		\sqrt{g(x)}\,\mathrm{d}z\,\mathrm{d}\xi\,\mathrm{d}x
		\;+\; \langle R_\infty\Phi_\varepsilon,\Phi_\varepsilon\rangle_{L^2}.
	\end{equation}
	Since $R_\infty$ has smooth kernel $K_{R_\infty}$, the remainder satisfies
	$|\langle R_\infty\Phi_\varepsilon,\Phi_\varepsilon\rangle|
	\le \|K_{R_\infty}\|_{L^\infty}\|\Phi_\varepsilon\|_{L^1}^2$.
	Changing variables $x = \exp_p(\varepsilon X)$ and using
	$\Phi_\varepsilon(\exp_p(\varepsilon X)) = \varepsilon^{-(n-1)/2}\Phi(X)$
	together with $\mathrm{d}v_{g_0} = \varepsilon^n\sqrt{g(\varepsilon X)}\,\mathrm{d}X$
	gives
	\[
	\|\Phi_\varepsilon\|_{L^1(\mathbb{S}^n)}
	= \varepsilon^{(n+1)/2}\bigl(\|\Phi\|_{L^1} + O_R(\varepsilon^2)\bigr),
	\]
	so $\|\Phi_\varepsilon\|_{L^1}^2 = O_R(\varepsilon^{n+1}) \to 0$, and thus
	the smoothing remainder is $o(1)$.
	
	In the main integral of \eqref{eq:osc}, substitute $x = \varepsilon X$,
	$z = \varepsilon Z$, $\xi = \varepsilon^{-1}\eta$, so
	$(x-z)\cdot\xi = (X-Z)\cdot\eta$ and
	$\mathrm{d}x\,\mathrm{d}z\,\mathrm{d}\xi = \varepsilon^n\,\mathrm{d}X\,\mathrm{d}Z\,\mathrm{d}\eta$.
	The two factors $\varepsilon^{-(n-1)/2}$ from $\Phi_\varepsilon(\varepsilon X)$
	and $\Phi_\varepsilon(\varepsilon Z)$, combined with $\varepsilon^n$, give the
	overall prefactor $\varepsilon^{-(n-1)/2}\cdot\varepsilon^{-(n-1)/2}\cdot
	\varepsilon^n = \varepsilon$. Hence
	\begin{equation}\label{eq:rescaled}
		\langle T\Phi_\varepsilon,\Phi_\varepsilon\rangle_{L^2}
		= (2\pi)^{-n}
		\iiint
		e^{i(X-Z)\cdot\eta}\,
		\bigl\langle
		\varepsilon\,a(\varepsilon X,\varepsilon^{-1}\eta)\,\Phi(Z),
		\Phi(X)
		\bigr\rangle
		\sqrt{g(\varepsilon X)}\,\mathrm{d}Z\,\mathrm{d}\eta\,\mathrm{d}X
		+ o(1).
	\end{equation}
	
	The density factor $\sqrt{g(\varepsilon X)}$ is a smooth $S^0$-multiplier
	satisfying $\sqrt{g(\varepsilon X)} = 1 + O_R(\varepsilon^2)$ uniformly for
	$|X| \le R$. Multiplying the rescaled symbols below by this factor therefore
	preserves all symbol-seminorm bounds, up to an additional $O_R(\varepsilon^2)$
	contribution that is subsumed in the $o_R(1)$ error.
	
	\smallskip
	\noindent\emph{(I) Lower-order terms.}
	For each \(j=0,-1,\ldots,-N_0\), put
	\[
	b_{\varepsilon,j}(X,\eta)
	:=
	\varepsilon\,a_j(\varepsilon X,\varepsilon^{-1}\eta).
	\]
	We split \(b_{\varepsilon,j}\) into low and high frequencies. Let
	\(\varphi\in C_c^\infty(\mathbb R^n)\) satisfy \(\varphi\equiv1\) near
	\(\eta=0\).
	
	For the low-frequency part we do not use an \(L^2\)-operator norm estimate.
	Since \(\Phi\) is fixed and compactly supported, \(\varphi b_{\varepsilon,j}\)
	is integrated only over compact \(X,Z,\eta\)-sets. On the support of \(\Phi\) and \(\varphi\), we have \(|X|\le R\) and
	\(\eta\in \operatorname{supp}\varphi\). Since \(a_j\in S^j\) with \(j\le0\),
	\[
	\|a_j(\varepsilon X,\varepsilon^{-1}\eta)\|_{\operatorname{End}(\mathbb C^N)}
	\le C_R.
	\]
	Hence
	\[
	\|b_{\varepsilon,j}(X,\eta)\|_{\operatorname{End}(\mathbb C^N)}
	=
	\varepsilon
	\|a_j(\varepsilon X,\varepsilon^{-1}\eta)\|_{\operatorname{End}(\mathbb C^N)}
	\le C_R\varepsilon.
	\]
	Therefore, using \(|\langle Av,w\rangle|\le \|A\|\,|v|\,|w|\) and integration by separation of variables,
	\[
	\begin{aligned}
		\left|
		\left\langle
		\operatorname{Op}(\varphi b_{\varepsilon,j})\Phi,\Phi
		\right\rangle
		\right|
		&\le
		C
		\iiint
		|\varphi(\eta)|
		\|b_{\varepsilon,j}(X,\eta)\|
		|\Phi(Z)|\,|\Phi(X)|
		\,dZ\,d\eta\,dX  \\
		&\le
		C_R\varepsilon
		\|\varphi\|_{L^1}
		\|\Phi\|_{L^1}^2
		=o_R(1).
	\end{aligned}
	\]
	
	For the high-frequency part \((1-\varphi)b_{\varepsilon,j}\), the estimates are
	uniform away from \(\eta=0\). For all multi-indices \(\alpha,\beta\),
	\[
	\partial_X^\alpha\partial_\eta^\beta b_{\varepsilon,j}(X,\eta)
	=
	\varepsilon^{1+|\alpha|-|\beta|}
	(\partial_x^\alpha\partial_\xi^\beta a_j)(\varepsilon X,\varepsilon^{-1}\eta).
	\]
	On \(\operatorname{supp}(1-\varphi)\), \(|\eta|\ge c>0\), and therefore
	\[
	|\partial_X^\alpha\partial_\eta^\beta b_{\varepsilon,j}(X,\eta)|
	\le
	C_{\alpha\beta R}\varepsilon^{1-j}(1+|\eta|)^{j-|\beta|}.
	\]
	Thus
	\[
	(1-\varphi)b_{\varepsilon,j}\in \varepsilon^{1-j}S^j
	\subset \varepsilon S^0 .
	\]
	Combining the low-frequency quadratic-form estimate with the high-frequency
	operator estimate, we obtain
	\[
	\bigl|\langle\operatorname{Op}(b_{\varepsilon,j})\Phi,\Phi\rangle\bigr|
	=o_R(1).
	\]
	Since the number of indices \(j=0,-1,\ldots,-N_0\) is finite, summing over
	these terms still gives an \(o_R(1)\) contribution to the quadratic form.	
	
	It remains to treat the finite-order remainder. Define
	\[
	b_{\varepsilon,\mathrm{rem}}(X,\eta)
	:=
	\varepsilon\,r_{-N_0-1}(\varepsilon X,\varepsilon^{-1}\eta).
	\]
	We split it again into low and high frequencies. The low-frequency contribution
	is estimated directly as above:
	\[
	\left|
	\left\langle
	\operatorname{Op}(\varphi b_{\varepsilon,\mathrm{rem}})\Phi,\Phi
	\right\rangle
	\right|
	=o_R(1),
	\]
	because \(\varphi b_{\varepsilon,\mathrm{rem}}\) is integrated only over compact
	\(X,Z,\eta\)-sets and satisfies
	\[
	\|\varphi(\eta)b_{\varepsilon,\mathrm{rem}}(X,\eta)\|_{End(\mathbb{C}^N)}
	\le C_R\varepsilon .
	\]
	For the high-frequency part, since \(r_{-N_0-1}\in S^{-N_0-1}\), the same
	rescaling calculation gives
	\[
	(1-\varphi)b_{\varepsilon,\mathrm{rem}}
	\in
	\varepsilon^{N_0+2}S^{-N_0-1}
	\subset	o_R(1)S^0.
	\]
	Thus its \(L^2\to L^2\) operator norm is \(o_R(1)\), and hence
	\[
	\left|
	\left\langle
	\operatorname{Op}((1-\varphi)b_{\varepsilon,\mathrm{rem}})\Phi,\Phi
	\right\rangle
	\right|
	=o_R(1).
	\]
	\smallskip
	\noindent\emph{(II) Principal symbol error.}
	Define
	\[
	r_\varepsilon(X,\eta) := \sqrt{g(\varepsilon X)}\,|\eta|_{g(\varepsilon X)} - |\eta|.
	\]
	We decompose $|\eta| = (1-\varphi(\eta))|\eta| + \varphi(\eta)|\eta|$, where
	$\varphi \in C_c^\infty(\mathbb{R}^n)$ satisfies $\varphi \equiv 1$ near
	$\eta = 0$. The high-frequency part $(1-\varphi)|\eta|$ is a classical symbol
	of order one, smooth away from the origin and supported in $|\eta| \ge c > 0$.
	
	The low-frequency part is treated directly at the level of quadratic forms.
	On the compact \(X\)-support of \(\Phi\) and on \(\operatorname{supp}\varphi\),
	the normal-coordinate expansion gives
	\[
	g^{ij}(\varepsilon X)=\delta^{ij}+O_R(\varepsilon^2),
	\qquad
	\sqrt{g(\varepsilon X)}=1+O_R(\varepsilon^2).
	\]
	Hence, on $\operatorname{supp}\varphi$,
	\[
	\left|	\sqrt{g(\varepsilon X)}\,|\eta|_{g(\varepsilon X)}-|\eta|
	\right|\leq \left| {(\sqrt {g(\varepsilon X)}  - 1)|\eta {|_{g(\varepsilon X)}}} \right| + \left| {|\eta {|_{g(\varepsilon X)}} - |\eta |} \right|	\le C_R\varepsilon^2|\eta|
	\le C_R\varepsilon^2.
	\]
	Therefore
	\[
	\begin{aligned}
		&\left|	(2\pi)^{-n}	\iiint
		e^{i(X-Z)\cdot\eta}
		\varphi(\eta)	r_\varepsilon(X,\eta)
		\langle \Phi(Z),\Phi(X)\rangle
		\,dZ\,d\eta\,dX	\right| \\
		&\qquad\le	C_R\varepsilon^2
		\|\varphi\|_{L^1}\|\Phi\|_{L^1}^2
		=o_R(1).
	\end{aligned}
	\] 
	The low-frequency contribution has therefore been controlled directly at the
	quadratic-form level. In the remaining high-frequency region, \(|\eta|\) is a
	smooth classical symbol. 
	
	For the
	high-frequency part, $g_{ij}(\varepsilon X) = \delta_{ij} + O_R(\varepsilon^2)$
	uniformly for $|X| \le R$, and the same chain-rule argument shows that the high-frequency error symbol \((1-\varphi)r_\varepsilon\) belongs to
	\(\varepsilon^2 S^1\) on the compact \(X\)-support. Therefore
	\[
	\operatorname{Op}((1-\varphi)r_\varepsilon):
	H^{1/2}(\mathbb R^n;\mathbb C^N)
	\to
	H^{-1/2}(\mathbb R^n;\mathbb C^N)
	\]
	has operator norm \(O_R(\varepsilon^2)\), and since
	\(\Phi\in C_c^\infty\subset H^{1/2}\),
	\[
	\bigl|	\langle	\operatorname{Op}((1-\varphi)r_\varepsilon)\Phi,\Phi	\rangle	\bigr|	\le	C_R\varepsilon^2\|\Phi\|_{H^{1/2}}^2
	=	O_R(\varepsilon^2).
	\]
	Collecting all contributions:
	\begin{equation}\label{eq:combined}
		\langle |\slashiii{D}|\Phi_\varepsilon,\Phi_\varepsilon\rangle_{L^2(\mathbb{S}^n)}
		= (2\pi)^{-n}
		\iiint
		e^{i(X-Z)\cdot\eta}\,|\eta|\,
		\bigl\langle\Phi(Z),\Phi(X)\bigr\rangle
		\,\mathrm{d}Z\,\mathrm{d}\eta\,\mathrm{d}X
		+ o_R(1).
	\end{equation}
	
	The triple integral in \eqref{eq:combined} is interpreted as the quadratic
	form of the Fourier multiplier $\operatorname{Op}(|\eta|)$. The Fourier multiplier \(\operatorname{Op}(|\eta|)\) is understood in the
	standard quadratic-form sense; the preceding cutoff decomposition justifies
	passing from the smooth symbol pieces to this multiplier. Since
	$\Phi \in C_c^\infty$, we have $\hat\Phi \in \mathcal{S}(\mathbb{R}^n;\mathbb{C}^N)$,
	and Plancherel's theorem gives
	\[
	(2\pi)^{-n}
	\iiint
	e^{i(X-Z)\cdot\eta}\,|\eta|\,
	\langle\Phi(Z),\Phi(X)\rangle
	\,\mathrm{d}Z\,\mathrm{d}\eta\,\mathrm{d}X
	= \langle\operatorname{Op}(|\eta|)\Phi,\Phi\rangle
	= (2\pi)^{-n}\int_{\mathbb{R}^n}|\eta|\,|\hat\Phi(\eta)|^2\,\mathrm{d}\eta.
	\]
	This equals $\int_{\mathbb{R}^n}\langle(-\Delta)^{1/2}\Phi,\Phi\rangle\,\mathrm{d}X$.
	The identification is consistent with the flat-space Dirac operator
	$\slashiii{D}_{\mathbb{R}^n}$, which satisfies
	\[
	\slashiii{D}_{\mathbb{R}^n}^2 = -\Delta\,\mathrm{Id}_{\mathbb{C}^N},
	\qquad\text{hence}\qquad
	|\slashiii{D}_{\mathbb{R}^n}|
	= \bigl(\slashiii{D}_{\mathbb{R}^n}^2\bigr)^{1/2}
	= (-\Delta)^{1/2}\,\mathrm{Id}_{\mathbb{C}^N}.
	\]
	Inserting into \eqref{eq:combined} and letting $\varepsilon \to 0$ completes
	the proof.
\end{proof}

\end{document}
